\chapter*{Abstract}
\addcontentsline{toc}{chapter}{Abstract}

Computer vision analysis in aquatic environments presents unique challenges due to light refraction, surface turbulence, and occlusion, which hinder the accuracy of standard motion tracking systems. This project presents a non-intrusive, AI-based proof-of-concept system for extracting swimming kinematics from uncalibrated, handheld poolside video footage. The core contribution is a novel methodology for dynamic pixel-to-meter scaling using custom 3D-printed underwater markers as calibration objects, addressing the limitations of fixed-camera setups across varying pool environments. A Model-in-the-Loop annotation pipeline, combining SAM 2 temporal segmentation with statistical outlier filtering and human verification, was developed to construct a domain-specific training dataset for a fine-tuned YOLO26 marker detector. Ego-motion compensation, ViTPose-based hip localisation, and Kalman filter smoothing were combined to produce absolute positional and velocity profiles from raw tracking output. Evaluation on a 50 m breaststroke sequence yielded a mean absolute split-time error of 1.005 s and an RMSE of 1.210 s across 18 checkpoints, demonstrating the feasibility of metric-accurate swimmer tracking without specialised hardware. Wrist keypoint tracking remained noisy in the underwater environment, highlighting distal pose estimation as the primary target for future refinement.